\chapter{Thiết Kế và Hiện Thực Hệ Thống}
\label{chap:system}

Chương này mô tả kiến trúc và hiện thực của \system{} -- từ tầng dữ liệu, engine năng lực, tích hợp KT thời gian thực, engine gợi ý, bài kiểm tra thích nghi, hệ thống chấm điểm, đến tầng phục vụ web và giao diện. Đây là phần kỹ thuật chuyển hóa các mô hình ở Chương~\ref{chap:theory} và~\ref{chap:experiment} thành một ứng dụng hoàn chỉnh chạy được.

% ─────────────────────────────────────────────────────────────────────────────
\section{Tổng Quan Kiến Trúc}
\label{sec:arch}

Hệ thống được tổ chức thành ba tầng (Hình~\ref{fig:arch}): (1) \textbf{tầng dữ liệu} chuẩn hóa question bank và taxonomy; (2) \textbf{tầng KT/scoring} huấn luyện offline và dự đoán năng lực; (3) \textbf{tầng phục vụ} cầu nối Python với giao diện React.

\begin{figure}[H]
  \centering
  \begin{tikzpicture}[
    node distance=0.55cm and 1.1cm,
    box/.style={draw, rounded corners, align=center, minimum height=0.9cm, minimum width=3.2cm, font=\small},
    layer/.style={draw, dashed, rounded corners, inner sep=0.35cm},
    arr/.style={-{Latex[length=2mm]}, thick},
  ]
    \node[box, fill=blue!8]   (qbank)  {Question Bank\\\scriptsize JSON (2.122 câu)};
    \node[box, fill=blue!8, right=of qbank] (schema) {Schemas + Taxonomy\\\scriptsize 17 KC};
    \node[box, fill=green!10, below=of qbank] (loader) {\texttt{data\_loader}\\\scriptsize normalize snake\_case};
    \node[box, fill=green!10, right=of loader] (comp) {\texttt{competency\_engine}\\\scriptsize SKILL\_GROUP\_TO\_KC};
    \node[box, fill=orange!12, below=of loader] (kt) {\texttt{kt\_predictor}\\\scriptsize DKT-Quality};
    \node[box, fill=orange!12, right=of kt] (rec) {\texttt{recommender}\\\scriptsize ZPD scoring};
    \node[box, fill=red!10, below=of kt] (fastapi) {FastAPI\\\scriptsize 9 endpoints \texttt{/api/*}};
    \node[box, fill=red!10, right=of fastapi] (grader) {\texttt{grader}\\\scriptsize Gemini + local};
    \node[box, fill=gray!12, below=of fastapi] (stream) {Streamlit\\\scriptsize iframe host};
    \node[box, fill=gray!12, right=of stream] (react) {React SPA\\\scriptsize JSX (no bundler)};

    \begin{scope}[on background layer]
      \node[layer, fit=(qbank)(schema), label=left:{\scriptsize\bfseries Dữ liệu}] {};
      \node[layer, fit=(loader)(comp), label=left:{\scriptsize\bfseries }] {};
      \node[layer, fit=(kt)(rec), label=left:{\scriptsize\bfseries KT/Scoring}] {};
      \node[layer, fit=(fastapi)(grader)(stream)(react), label=left:{\scriptsize\bfseries Phục vụ}] {};
    \end{scope}

    \draw[arr] (qbank) -- (loader);
    \draw[arr] (schema) -- (comp);
    \draw[arr] (loader) -- (kt);
    \draw[arr] (comp) -- (rec);
    \draw[arr] (kt) -- (fastapi);
    \draw[arr] (rec) -- (grader);
    \draw[arr] (fastapi) -- (stream);
    \draw[arr] (grader) -- (react);
    \draw[arr] (react.east) to[bend right=35] node[right,font=\scriptsize]{fetch \texttt{/api}} (fastapi.east);
  \end{tikzpicture}
  \caption{Kiến trúc ba tầng của \system{}: dữ liệu $\to$ KT/scoring $\to$ phục vụ.}
  \label{fig:arch}
\end{figure}

Một đặc điểm quan trọng là backend Python và frontend React dùng \textbf{hai quy ước đặt tên khác nhau} (snake\_case vs camelCase) và hai lược đồ câu hỏi khác nhau (\texttt{skill\_groups} vs \texttt{skillGroups}, \texttt{question\_id} vs \texttt{id}). Việc chuẩn hóa xảy ra ở \textit{ranh giới} (boundary) tại hai nơi: \texttt{\_map\_question()} (Python $\to$ React) và \texttt{\_normalize\_question()} / \texttt{\_normalize\_candidate\_response()} (React $\to$ Python). Thiết kế này giữ mỗi phía sạch theo quy ước riêng và tập trung mọi phép chuyển đổi vào một chỗ.

% ─────────────────────────────────────────────────────────────────────────────
\section{Tầng Dữ Liệu và Engine Năng Lực}
\label{sec:data_layer}

\subsection{Chuẩn hóa câu hỏi}
\texttt{data\_loader.py} nạp \texttt{question\_bank.json} và chuẩn hóa về schema Python (snake\_case). Mỗi câu hỏi có các trường: \texttt{question\_id}, \texttt{question\_text}, \texttt{roles.primary}, \texttt{difficulty\_label}, \texttt{difficulty\_score}, \texttt{question\_type}, \texttt{skill\_groups}, \texttt{skill\_tags}, \texttt{answers} (gồm \texttt{detailed}, \texttt{evaluation\_points}, \texttt{explanation}, $\ldots$).

\subsection{Ánh xạ KC (Competency Engine)}
\texttt{competency\_engine.py} định nghĩa taxonomy 17 KC theo vai trò và bảng tra \texttt{SKILL\_GROUP\_TO\_KC} (82 ánh xạ) gộp các tag chi tiết về KC chuẩn. Ví dụ tất cả \texttt{SQL\_WINDOW\_FUNCTION}, \texttt{SQL\_JOIN}, \texttt{SQL\_CTE}, \texttt{DATABASE\_DESIGN}$\ldots$ đều ánh xạ về KC \texttt{SQL\_DATABASE}. Engine cũng cung cấp hai ngưỡng phân loại năng lực:
\begin{equation}
  \text{KC yếu nếu } \theta < 0.40; \qquad \text{KC mạnh nếu } \theta \geq 0.70 .
\end{equation}
cùng các hàm \texttt{identify\_weak\_kcs()}, \texttt{identify\_strong\_kcs()}, \texttt{overall\_score()}.

% ─────────────────────────────────────────────────────────────────────────────
\section{Tích Hợp Knowledge Tracing Thời Gian Thực}
\label{sec:kt_runtime}

Khi chạy, \texttt{kt\_predictor.py} nạp mô hình tốt nhất theo ablation (\texttt{dkt\_quality} mặc định) một lần (singleton) và dự đoán mức thành thạo $\in [0.05, 0.99]$ cho mỗi KC từ lịch sử tương tác. Nếu lịch sử quá ngắn ($< H_{\min}=3$) hoặc không nạp được mô hình (môi trường không có PyTorch), predictor trả về \texttt{\{\}} và hệ thống tự động lui về EMA thuần -- cơ chế \textit{graceful fallback}.

Skill vector được cập nhật theo Công thức~\eqref{eq:blend}: EMA luôn chạy trước, sau đó blend KT(70\%) + EMA(30\%) khi đủ lịch sử. Hằng số 70/30 và ngưỡng 3 tương tác được giữ \textbf{đồng nhất} ở backend (\texttt{recommendation\_engine.\_kt\_predict}) và ở frontend (hàm \texttt{mstSkillVector}) để con số hiển thị trên Dashboard và màn KT Model luôn khớp nhau.

% ─────────────────────────────────────────────────────────────────────────────
\section{Engine Gợi Ý theo Zone of Proximal Development}
\label{sec:recsys_impl}

\texttt{recommender.py} xếp hạng câu hỏi theo điểm tổng hợp. Khi tập trung vào KC yếu (\texttt{focus\_weak=True}):
\begin{equation}
  \text{score}(q) = 0.65 \cdot \text{weakness} + 0.30 \cdot \text{diff\_fit} + 0.05 \cdot \text{coverage},
  \label{eq:rec_weak}
\end{equation}
trong đó \textit{weakness} $= \max(0, 0.40 - \theta_{kc})/0.40$ ưu tiên KC dưới ngưỡng yếu; \textit{diff\_fit} đo độ phù hợp giữa độ khó câu hỏi và năng lực hiện tại (chiến lược ZPD -- chọn câu hơi khó hơn năng lực); \textit{coverage} thưởng cho việc phủ KC chưa mạnh. Khi ở chế độ thử thách (\texttt{focus\_weak=False}), trọng số đổi thành $0.45$ readiness $+ 0.45$ diff\_fit $+ 0.10$ coverage.

\texttt{recommendation\_engine.py} điều phối: gọi KT predictor $\to$ blend EMA $\to$ \texttt{recommend\_questions()}, đồng thời cung cấp các API mức cao: \texttt{get\_next\_questions()}, \texttt{get\_weak\_kc\_practice()}, \texttt{get\_challenge\_questions()}, \texttt{get\_session\_summary()}.

% ─────────────────────────────────────────────────────────────────────────────
\section{Bài Kiểm Tra Thích Nghi (Multi-Stage Testing)}
\label{sec:mst}

Phiên đánh giá (Diagnostic) áp dụng CAT hai giai đoạn (Hình~\ref{fig:mst}):

\begin{figure}[H]
  \centering
  \begin{tikzpicture}[
    node distance=0.9cm and 1.4cm,
    b/.style={draw, rounded corners, align=center, font=\small, minimum height=0.85cm},
    arr/.style={-{Latex[length=2mm]}, thick}, font=\small]
    \node[b, fill=blue!8] (s1) {Stage 1 -- Placement\\\scriptsize câu phủ đều các KC};
    \node[b, fill=gray!12, right=2.2cm of s1] (route) {Routing\\\scriptsize theo điểm TB};
    \node[b, fill=red!10, above right=0.2cm and 1.6cm of route] (weak) {WEAK};
    \node[b, fill=orange!12, right=1.6cm of route] (mid) {MID};
    \node[b, fill=green!10, below right=0.2cm and 1.6cm of route] (strong) {STRONG};
    \node[b, fill=blue!6, right=5.6cm of route] (s2) {Stage 2 -- Adaptive\\\scriptsize ưu tiên KC yếu};
    \draw[arr] (s1) -- (route);
    \draw[arr] (route) -- (weak);
    \draw[arr] (route) -- (mid);
    \draw[arr] (route) -- (strong);
    \draw[arr] (weak) -- (s2);
    \draw[arr] (mid) -- (s2);
    \draw[arr] (strong) -- (s2);
  \end{tikzpicture}
  \caption{Luồng Multi-Stage Testing: Stage 1 định vị $\to$ phân nhánh $\to$ Stage 2 thích nghi.}
  \label{fig:mst}
\end{figure}

\textbf{Stage 1 (định vị):} chọn câu phủ đều các KC của vai trò. Điểm trung bình quyết định \textit{routing}: WEAK / MID / STRONG. \textbf{Stage 2 (thích nghi):} chọn câu chuyên sâu, ưu tiên các KC yếu nhất phát hiện ở Stage 1, độ khó điều chỉnh theo nhánh (STRONG $\to$ ưu tiên ADVANCED; WEAK $\to$ INTERMEDIATE để học). Sau cùng, vector năng lực được cập nhật \textbf{theo từng trục KC}: chỉ trục được kiểm tra mới đổi giá trị; lần đánh giá đầu tiên ánh xạ trực tiếp điểm số $\to$ giá trị KC (vì không có baseline để cộng dồn delta).

Điểm cuối được tính bằng chính pipeline KT--EMA dùng chung (\texttt{mstSkillVector}), nên radar ở màn kết quả, Dashboard và màn KT Model đều hiển thị cùng một con số.

% ─────────────────────────────────────────────────────────────────────────────
\section{Hệ Thống Chấm Điểm}
\label{sec:grading}

\texttt{grader.py} chấm câu trả lời theo loại:
\begin{itemize}[leftmargin=1.3em]
  \item \textbf{Câu khách quan} (MC\_SINGLE, TRUE\_FALSE, FILL\_BLANK): chấm \textit{tất định} phía client/local, không gọi LLM.
  \item \textbf{Câu tự luận / coding} (THEORY, PRACTICE, CODING): chấm bằng \textbf{Gemini} (\texttt{gemini-3.5-flash}) với prompt chứa câu hỏi, đáp án tham chiếu và \texttt{evaluation\_points}.
\end{itemize}
Khi Gemini lỗi (rate limit, mạng), hệ thống lui về \textbf{local rubric} (\texttt{scoring.py}): mỗi \texttt{evaluation\_point} được coi là ``đạt'' khi $\geq 50\%$ từ khóa của nó xuất hiện trong câu trả lời; nếu không có eval\_points thì dùng độ tương đồng Jaccard từ khóa so với đáp án tham chiếu. Lỗi Gemini được trả kèm trong trường \texttt{gemini\_error} để dễ chẩn đoán.

% ─────────────────────────────────────────────────────────────────────────────
\section{Tầng Phục Vụ và Giao Diện}
\label{sec:serving}

\subsection{Vì sao tách FastAPI khỏi Streamlit}
Streamlit không phục vụ tốt tệp tĩnh với đúng MIME type (trả \texttt{text/plain} cho \texttt{.html}, trình duyệt từ chối render). Do đó UI thực \textbf{không} do Streamlit phục vụ: \texttt{app.py} khởi động một server FastAPI ở thread nền (cổng 8000), FastAPI mount thư mục \texttt{static/} qua \texttt{StaticFiles(html=True)} để uvicorn đặt đúng content-type; trang Streamlit chỉ là một \texttt{<iframe>} trỏ vào UI. React sau đó gọi ngược về \texttt{/api/$\ldots$} qua \texttt{fetch}.

\begin{table}[H]
  \centering
  \caption{Các endpoint FastAPI của \system{}}
  \label{tab:endpoints}
  \begin{tabular}{lll}
    \toprule
    \textbf{Method} & \textbf{Endpoint} & \textbf{Chức năng} \\
    \midrule
    GET  & \texttt{/api/health}                 & Kiểm tra trạng thái server \\
    POST & \texttt{/api/grade}                  & Chấm điểm một câu trả lời \\
    POST & \texttt{/api/user/create}            & Tạo người dùng mới \\
    GET  & \texttt{/api/user/\{id\}}            & Lấy hồ sơ người dùng \\
    GET  & \texttt{/api/user/\{id\}/summary}    & Tóm tắt tiến độ \\
    POST & \texttt{/api/user/\{id\}/answer}     & Ghi nhận câu trả lời, cập nhật skill \\
    GET  & \texttt{/api/recommend/\{id\}}       & Gợi ý câu hỏi tiếp theo \\
    POST & \texttt{/api/assessment/stage1}      & Sinh câu hỏi Stage 1 \\
    POST & \texttt{/api/assessment/finalize}    & Chốt đánh giá, chạy KT pipeline \\
    \bottomrule
  \end{tabular}
\end{table}

\subsection{Giao diện React không bundler}
Frontend gồm các tệp \texttt{.jsx} nạp tuần tự bằng thẻ \texttt{<script type="text/babel">} và biên dịch ngay trên trình duyệt bằng Babel standalone -- không cần npm/webpack. Thứ tự nạp cố định: \texttt{state $\to$ chrome $\to$ widgets $\to$ screens-* $\to$ router}. Vì không có module system, mọi tệp chia sẻ một global scope; \texttt{state.jsx} định nghĩa store toàn cục (React Context + sessionStorage), dữ liệu mock và sẽ được ghi đè bởi \texttt{window.\_\_INTERNHUB\_DATA\_\_} (question bank serialize từ Python). Đánh đổi: một lỗi cú pháp ở bất kỳ tệp nào sẽ làm hỏng toàn bộ app, nên cần kiểm thử parse cẩn thận.

% ─────────────────────────────────────────────────────────────────────────────
\section{Minh Họa Hệ Thống}
\label{sec:demo}

% (Ảnh minh họa giao diện có thể bổ sung vào figures/ nếu cần.)


Toàn bộ luồng -- onboarding chọn vai trò $\to$ Diagnostic (MST) $\to$ chấm điểm $\to$ cập nhật vector $\to$ Dashboard/Roadmap/KT Model -- hoạt động end-to-end trong một codebase chạy local.
